\chapter{Advanced Segmentation Architectures}
\label{ch:advseg}

\chapmeta{Advanced}{$\sim$5\,h (2\,h + 3\,h)}{LoveDA}

\begin{objbox}
\begin{itemize}[leftmargin=1.2em,nosep]
  \item Use \texttt{segmentation-models-pytorch} (SMP) to build decoders on pretrained encoders.
  \item Understand atrous (dilated) convolution and DeepLabV3+'s ASPP module.
  \item Understand multi-scale context: FPN and PSPNet pyramid pooling.
  \item Benchmark architectures fairly on a common dataset and budget.
  \item Choose an architecture/encoder for a given EO constraint.
\end{itemize}
\end{objbox}

The hand-built U-Net of Chapter~\ref{ch:seg} taught the mechanics. In practice
we assemble decoders on top of pretrained encoders with a library, and choose
among architectures designed to capture \emph{multi-scale context} --- essential
in EO, where land-cover objects range from single trees to whole agricultural
regions.

\section{segmentation-models-pytorch}

SMP exposes every architecture through one uniform constructor, pairing a
decoder with any of dozens of pretrained encoders (\texttt{timm}-backed):
\begin{lstlisting}
import segmentation_models_pytorch as smp
model = smp.Unet(encoder_name="resnet50", encoder_weights="imagenet",
                 in_channels=3, classes=7)
# swap decoder by changing the class: smp.DeepLabV3Plus / smp.FPN / smp.PSPNet
\end{lstlisting}
For multispectral input set \texttt{in\_channels=13}; SMP adapts the first
convolution automatically (cf.\ Chapter~\ref{ch:transfer}). This one-line
flexibility is what makes principled architecture comparison feasible.

\section{Atrous convolution and DeepLabV3+}

A standard convolution must pool to enlarge its receptive field, sacrificing
resolution. \emph{Atrous} (dilated) convolution inserts gaps of rate $r$ between
kernel taps, enlarging the receptive field \emph{without} downsampling or extra
parameters. A dilated $k\times k$ kernel has an effective size
\begin{equation}
k_\mathrm{eff} = k + (k-1)(r-1),
\end{equation}
so a $3\times3$ kernel at rate $r{=}6$ covers $13\times13$. DeepLabV3+ uses
\emph{Atrous Spatial Pyramid Pooling} (ASPP): several parallel atrous
convolutions at different rates (e.g.\ $6,12,18$) plus image-level global
pooling, concatenated to fuse context at multiple scales, followed by a light
decoder that recovers boundaries. It excels when objects span very different
sizes --- the norm in land cover.

\section{FPN and PSPNet}

\emph{Feature Pyramid Networks} (FPN) build a top-down pathway with lateral
connections, producing semantically strong features at every resolution and
fusing predictions across the pyramid --- naturally multi-scale and efficient.
\emph{PSPNet} introduces the \emph{pyramid pooling module}, pooling the deepest
feature map at several grid scales ($1\times1, 2\times2, 3\times3, 6\times6$),
processing each, and concatenating, injecting global scene context that helps
disambiguate locally-similar classes (e.g.\ shadowed roof vs water). All three
--- DeepLabV3+, FPN, PSPNet --- attack the same problem (context at multiple
scales) by different means.

\section{A fair benchmark}

Comparisons are only meaningful under a fixed protocol: same encoder, same data
splits, same augmentation, same epochs and optimiser, varying only the decoder.
Report mIoU, parameters, and inference latency.

\begin{center}\small
\renewcommand{\arraystretch}{1.2}
\begin{tabular}{@{}lllll@{}}
\toprule
Decoder & Context mechanism & Params & Strength & Typical use\\
\midrule
U-Net & skip connections & low--med & sharp boundaries & default, small data\\
FPN & feature pyramid & low & multi-scale, fast & efficient inference\\
PSPNet & pyramid pooling & med & global context & scene-level classes\\
DeepLabV3+ & ASPP (atrous) & med--high & multi-scale + edges & best all-round\\
\bottomrule
\end{tabular}
\end{center}

\noindent On LoveDA with a ResNet50 encoder the four are usually within a few
mIoU points; DeepLabV3+ tends to lead on boundary-heavy urban scenes, U-Net
remains a strong, low-memory default, and FPN is attractive when latency
matters. The lesson is to \emph{measure on your data and budget} rather than
trust leaderboard ranks from another domain.

\begin{keybox}
\begin{itemize}[leftmargin=1.2em,nosep]
  \item SMP $=$ decoder $\times$ pretrained encoder via one constructor; trivial multispectral support.
  \item Atrous convolution grows the receptive field without losing resolution; $k_\mathrm{eff}=k+(k-1)(r-1)$.
  \item DeepLabV3+ (ASPP), FPN and PSPNet all capture multi-scale context differently.
  \item Benchmark fairly: fix everything but the decoder; report mIoU, params, latency.
  \item No universal winner --- choose for your data, objects' scale, and compute budget.
\end{itemize}
\end{keybox}

\begin{practbox}
Notebook: \nbpath{chapter\_10\_advanced\_segmentation/10\_practical\_advanced\_segmentation.ipynb}.
\begin{enumerate}[leftmargin=1.4em,nosep]
  \item Build U-Net, DeepLabV3+, FPN and PSPNet with SMP on a shared ResNet50
  encoder; print parameter counts.
  \item Visualise how atrous rate enlarges the receptive field.
  \item Run a fixed-protocol benchmark of the four decoders for a few epochs.
  \item Plot mIoU vs params vs latency; pick a winner.
  \item Train the winner longer and display prediction triplets.
\end{enumerate}
\end{practbox}

\begin{exbox}
\textbf{10.1 Encoder swap.} Hold the decoder fixed and compare ResNet34/50 and
EfficientNet-B0/B4 encoders by mIoU and speed.\\[2pt]
\textbf{10.2 ASPP rates.} Vary DeepLab's atrous rates; relate them to the object
scales in LoveDA.\\[2pt]
\textbf{10.3 Multispectral.} Rebuild the best model with \texttt{in\_channels=13}
on a multispectral dataset; measure the gain over RGB.\\[2pt]
\textbf{Mini-project.} Produce a one-page ``architecture selection'' report
recommending a decoder+encoder for (a) a latency-constrained edge device and (b)
a maximum-accuracy server, justified by your benchmark.
\end{exbox}
